\documentclass{article}

% "preprint" option is used for arXiv or other preprint submissions
\usepackage[preprint]{neurips_2025}

\usepackage[utf8]{inputenc} % allow utf-8 input
\usepackage[T1]{fontenc}    % use 8-bit T1 fonts
\usepackage{hyperref}       % hyperlinks
\usepackage{url}            % simple URL typesetting
\usepackage{booktabs}       % professional-quality tables
\usepackage{amsfonts}       % blackboard math symbols
\usepackage{dsfont}         % extra math symbols (indicator function)
\usepackage{nicefrac}       % compact symbols for 1/2, etc.
\usepackage{microtype}      % microtypography
\usepackage{xcolor}         % colors
\usepackage{graphicx}       % pictures

% Additional packages (personal preference)
\usepackage{amsmath, amsfonts}
\usepackage{subcaption}
\usepackage{multirow}

% Note. For the workshop paper template, both \title{} and \workshoptitle{} are required, with the former indicating the paper title shown in the title and the latter indicating the workshop title displayed in the footnote.
\title{ \Large{ICCAD 2026 CAD Contest Problem D} \\ \LARGE{Timing Fixing by Useful Skew}}

% The \author macro works with any number of authors. There are two commands
% used to separate the names and addresses of multiple authors: \And and \AND.
%
% Using \And between authors leaves it to LaTeX to determine where to break the
% lines. Using \AND forces a line break at that point. So, if LaTeX puts 3 of 4
% authors names on the first line, and the last on the second line, try using
% \AND instead of \And before the third author name.

\author{%
  Aaron~Wu \\
  Department of Electrical Engineering\\
  National Taiwan University\\
  \texttt{b13901011@ntu.edu.tw} \\
  \And
  Yucheng~Liu \\
  Department of Electrical Engineering\\
  National Taiwan University\\
  \texttt{b13901104@ntu.edu.tw} \\
  \And
  Sheng-Yu~Cheng \\
  Department of Electrical Engineering\\
  National Taiwan University\\
  \texttt{b13901088@ntu.edu.tw} \\
  \And
  Hong-Bin~Lai \\
  Department of Electrical Engineering\\
  National Taiwan University\\
  \texttt{b13901078@ntu.edu.tw}
}

\begin{document}
%%%%%%%%%
% for consistency, We have to unify the synonyms as follows:
% agent = model = policy
%%%%%
\maketitle

\begin{abstract}
  % Briefly summarize the problem, the useful-skew optimization approach, the
  % implemented candidate/acceptance policy framework, and the main experimental
  % finding. Keep this to one compact paragraph for a 6--8 page report.
\end{abstract}

\section{Introduction}
ICCAD Contest 2026 Problem D asks us to improve such violations by modifying the clock tree rather than changing the data-path logic.  The available operations are local clock-tree edits: inserting a new buffer, removing an inserted buffer, or resizing an existing buffer.  These operations change clock arrival times at flip-flops and therefore change the skew between launch and capture registers.

The key opportunity is useful skew.  A nonzero clock skew is not always harmful:
delaying the capture clock can give a setup-critical data path more time, while
delaying the launch clock can help avoid hold violations by postponing data
launch.  However, each clock-tree edit may affect many data paths at once.
Improving one endpoint may degrade another path, and adding buffers also
increases area.  The optimization problem is therefore not only to remove timing
violations, but to balance setup improvement, hold improvement, and area cost.

In this project, we model the clock tree as a mutable state and search over a
space of legal buffer insertion, removal, and resizing operations.  Our solver
uses candidate policies to decide which moves should be evaluated, and
acceptance policies to decide which evaluated move should be applied.  This
framework allows us to compare greedy search, simulated annealing, iterated
simulated annealing, two-step timing repair, and tabu search under a common
clock-tree timing model.  The experimental discussion focuses on how these
policies trade off setup slack, hold slack, and area.

\section{Problem Formulation}
\subsection{Slack, Skew, and Score}
A data path connects a launch flip-flop to a capture flip-flop through
combinational logic.  Let $D$ be the data-path delay, $T_{\mathrm{clk}}$ be the
clock period, $T_{\mathrm{setup}}$ be the setup requirement, and
$T_{\mathrm{hold}}$ be the hold requirement.  Let
$C_{\mathrm{launch}}$ and $C_{\mathrm{capture}}$ denote the clock arrival times
from the clock source to the launch and capture flip-flops.  We define clock
skew as
\begin{equation}
  \mathrm{skew} = C_{\mathrm{capture}} - C_{\mathrm{launch}}.
\end{equation}
With this convention, the setup and hold slacks used in the presentation are
\begin{align}
  \mathrm{setup\_slack}
    &= T_{\mathrm{clk}} - T_{\mathrm{setup}} - D + \mathrm{skew}, \\
  \mathrm{hold\_slack}
    &= D - T_{\mathrm{hold}} - \mathrm{skew}.
\end{align}
A negative setup slack means that data arrives too late at the capture
flip-flop.  A negative hold slack means that data arrives too early and may
overwrite the previous value.  Therefore, a setup violation can often be
improved by increasing the capture clock delay, while a hold violation can often
be improved by increasing the launch clock delay.

The objective combines timing and area.  Timing quality is measured by total
negative slack (TNS) and worst negative slack (WNS) under the setup-oriented SS
corner and the hold-oriented FF corner.  Area is the total area of the clock-tree
buffers.  Given a baseline metric value $x_0$ and a final value $x$, we use the
relative improvement form
\begin{equation}
  I(x, x_0) = 1 - \frac{x}{x_0},
\end{equation}
when $x_0$ is nonzero.  This rewards reducing negative timing metrics toward
zero and rewards reducing area below the baseline.  The implemented score is
\begin{align}
  S_{\mathrm{setup}}
    &= I(\mathrm{TNS}_{SS}, \mathrm{TNS}_{SS}^{0})
     + I(\mathrm{WNS}_{SS}, \mathrm{WNS}_{SS}^{0}), \\
  S_{\mathrm{hold}}
    &= I(\mathrm{TNS}_{FF}, \mathrm{TNS}_{FF}^{0})
     + I(\mathrm{WNS}_{FF}, \mathrm{WNS}_{FF}^{0}), \\
  S_{\mathrm{area}}
    &= 1 - \frac{\mathrm{Area}}{\mathrm{Area}^{0}}, \\
  S
    &= \alpha S_{\mathrm{setup}}
     + \beta S_{\mathrm{hold}}
     + \gamma S_{\mathrm{area}}. , \\
     \alpha = 0.5, \beta = 0.25, \gamma = 0.25
\end{align}
This formula explains the main trade-off in the problem: aggressive buffer
insertion may improve setup or hold slack, but the same move can reduce the area
component and may also hurt timing on other paths. The actual parameter is not given; therefore, we set parameters based on the importance for evaluation.

\subsection{Optimization Moves}
The state space consists of all clock trees reachable from the input tree by
legal edits.  A state records which buffers are present, the cell type of each
buffer, and the parent-child connectivity of the clock tree.  Original contest
nodes remain part of the design, while buffers inserted by the optimizer may be
removed later.

Our optimizer uses three move types.  First, an \emph{insert} move places a new
buffer on an existing clock-tree edge.  This increases the clock delay of the
downstream subtree and can be used to adjust useful skew, but it also increases
area.  Second, a \emph{remove} move deletes a buffer that was inserted by the
optimizer and reconnects the surrounding tree.  This can recover area, but it
may undo a previous timing repair.  Third, a \emph{resize} move changes the cell
type of an existing buffer.  Resizing gives a finer way to increase or decrease
clock delay and area without changing the clock-tree topology.

Because the number of possible edit sequences is large, the solver cannot
enumerate the full state space.  Instead, each iteration generates a smaller set
of candidate moves and evaluates their impact on the current timing and score.
The rest of the report describes how we construct these candidate sets and how
different search policies decide which move to accept.

\section{Methodology}
\subsection{Model for Optimization}
% Briefly describe ClockTree, DataPathGraph, and TimingState.
% Emphasize incremental timing evaluation and reversible apply/undo edits.

\textbf{Clock Tree.}
The clock tree is the primary data structure in our optimization framework. It consists of the clock source as the root, buffers as internal nodes, and flip-flops as leaves. The supported operations include buffer insertion, buffer removal, and buffer resizing.

\medskip

\textbf{Data Path Graph.}
The data path graph stores the launch flip-flop, capture flip-flop, and data-path delay information. This information is required to evaluate whether a clock-tree modification improves timing slack.

\medskip

\textbf{Timing State.}
The timing state maintains timing and area metrics, including setup slack, hold slack, total negative slack (TNS), worst negative slack (WNS), and area. These metrics are updated and queried during optimization.

\medskip

\textbf{Design Philosophy.}

\textit{Minimize Re-computation.}
Recomputing timing information for the entire clock tree after every modification is computationally expensive. Therefore, we employ an incremental timing cache that updates only the affected paths instead of the entire structure.

\medskip

\textit{Recoverability.}
The optimization process must support undo operations so that a modification can be reverted efficiently when the move is rejected.

\subsection{Candidate Policies}
% Group the policies instead of describing every optimizer in long form:
% random action-space sampling, violation-driven policies, upstream/critical
% endpoint policies, and union/sampled-union pools.
\textbf{Violation Path.}
We check all data paths and find out those whose slack is negative and smallest, which means the data path violate timing constrain. If setup constrain is violated, the buffers near capture flip-flop is modified. Otherwise, if hold constrain is violated, the buffers near launch flip-flop is modified

\medskip

\textbf{Upstream Window.}
We look for violation endpoint (i.e. capture flip-flop) and explore several buffers toward upstream of clock tree. The policy adds more candidate for optimization, but also more time consuming.

\medskip

\textbf{Critical Endpoint.}
We define \textit{criticality} as an evaluation metric for how much a flip-flop worsen total negative slack. Negative setup slack is added to violation path's capture flip-flop's criticality, and negative hold slack is added to launch flip-flop's criticality, respectively. We try to modify the flip-flop with largest criticality to optimize the violation on more than one data path.


\subsection{Accept Policies}
% Explain BestScore greedy, Metropolis/iterated SA, tabu search, and two-step
% slack-then-score optimization. A compact table can summarize optimizer name,
% candidate policy, accept policy, and key idea.
\textbf{Greedy Approach.}
Based on chosen candidate policy, we apply all generated candidate operations to current clock tree and calculate $\Delta Score$. The strategy accept only operations with $\Delta Score>0$ and the largest $\Delta Score$. 

\medskip

\textbf{Two Step Optimization.}
It's similar to greedy approach but with two stage focusing on different objective. In the first stage, we try to optimize timing constrain and neglect the penalty caused by area increment. In the second stage, we try to remove and resize buffers to reduce area penalty with timing performance maintained.

\medskip

\textbf{Simulated Annealing.}
The strategy randomly come up with an move and apply the move on current clock tree. If $\Delta Score>0$ then the move is accepted, otherwise with a certain probability determined by the temperature, this move will be accepted. When the temperature is low, the strategy act like greedy approach.
\textit{Iterated Simulated Annealing.}
The strategy take turns doing standard simulated annealing step and greedy step, trying to explore more possible moves.

\medskip

\textbf{Tabu Search.}
In each step of the strategy, we pick the best move among all candidates greedily. If the move is non-tabu move, then we directly accept the move, otherwise only if the move creates new global optimal will the move be accepted.

\section{Experimental Results and Discussion}
\subsection{Overall Comparison}
% Present the main ranking table or figure. Focus on the final score and the
% best-performing optimizer.

\subsection{Setup, Hold, and Area Trade-off}
% Break down the result by setup component, hold component, and area component.
% Highlight cases where timing improves but area penalty reduces final score.

\subsection{Policy-Level Discussion}
% Discuss what the candidate-policy and accept-policy comparisons imply.
% Example: random neighborhoods work well with tabu; sampled union pools improve
% timing aggressively but can increase area.

\subsection{Per-Testcase Behavior}
% Summarize testcase-dependent behavior only if space allows. Keep this short
% and focus on why no single policy dominates every testcase.

\section{Conclusion and Future Work}
% Summarize the key findings in a short paragraph. Mention one or two future
% directions, such as area-aware tabu polish or adaptive candidate selection.

\section*{Author Contributions}

\section*{Generative AI Usage Disclosure}
We used generative AI tools during the development of this project. In particular, OpenAI Codex and Cursor was used extensively to assist with source-code implementation, including parser construction, clock-tree data structures, timing evaluation routines, optimization-loop implementation, test scripts, and refactoring. Approximately 80\% of the submitted source code was initially generated or modified with the assistance of Codex or Cursor.

The algorithmic ideas, including the useful-skew optimization strategy, candidate generation policies, acceptance strategies, timing-area trade-off considerations, experimental design, and final optimizer selection, were designed and evaluated by the team members. All AI-generated code was manually reviewed, integrated, debugged, and tested by the team. The reported experimental results were obtained by executing our implemented solver on the test cases; no experimental data or performance numbers were generated by AI tools.

We also used ChatGPT/Codex to assist with report organization, wording refinement, and documentation polishing. The final technical content, claims, results, and conclusions were checked and approved by the team.


\medskip
{
  \small
  \bibliographystyle{unsrt}
  \bibliography{references}
}
%%%%%%%%%%%%%%%%%%%%%%%%%%%%%%%%%%%%%%%%%%%%%%%%%%%%%%%%%%%%

\end{document}
